\section{Conclusion}
\label{sec:conclusion}

This paper studied MLU-oriented traffic engineering for LEO satellite
constellations when telemetry constraints expose only a small fraction of
demands. We presented \sysname{}, a sparse-aware learned allocator that
represents missing measurements explicitly, prevents placeholder traffic
from contaminating graph messages, and combines spatial and temporal
evidence before predicting path split ratios. Shared parameters and demand
pruning keep the architecture independent of the instantaneous demand-set
size.

Repeated evaluations on a 1{,}584-satellite Starlink constellation show
that MLU rankings vary with both the observation ratio and model
initialization.
\sysname{} attains the lower MLU on more matched trials than the zero-fill,
mean-interp, and nbr-fill controls, yet none of its external pairwise
comparisons is statistically significant. The component diagnostic is also
mixed at the two sparsest ratios. These negative results matter: they show
that the structural plausibility of a sparse-aware representation does not
by itself establish a uniform end-to-end MLU gain, and that comparisons
must account for variation across experimental repeats rather than rely on
a single summary statistic.

For deployment, a source-neighbor estimate can provide optional
post-processing with a useful congestion signal without access to the
complete traffic matrix. It is a supporting mechanism, not a substitute for
the sparse-aware allocator. Future work should improve robustness across
initializations and observation masks through adaptive measurement
scheduling, observation-aware model selection, and objectives that
explicitly control worst-case MLU.
